\chapter*{}

\begin{center}
\textbf{\large Résumé : }
\end{center}

Ce rapport présente le développement de \textbf{SafeTrack}, une application web et mobile intelligente dédiée à la sécurité et au suivi en temps réel des enfants et des personnes vulnérables. Basée sur une architecture microservices, la plateforme intègre plusieurs fonctionnalités telles que l’authentification sécurisée, le suivi GPS en temps réel, la gestion des zones sécurisées (géofencing), les alertes SOS, les notifications instantanées et l’analyse prédictive des risques par intelligence artificielle. 

Le système exploite des techniques avancées d’IA, notamment la \textbf{Retrieval-Augmented Generation (RAG)}, afin de fournir des analyses contextuelles et des recommandations de sécurité adaptées aux situations détectées. L’application permet également la communication en temps réel entre parents, aidants et institutions pour renforcer la protection proactive des utilisateurs.

Les principales technologies utilisées sont : \textbf{Flutter}, \textbf{NestJS}, \textbf{Python}, \textbf{MongoDB}, \textbf{Redis}, \textbf{WebSocket}, ainsi que des modèles d’intelligence artificielle intégrés via \textbf{Ollama} et \textbf{LangChain}.\\

\noindent
\textbf{Mots clés :} SafeTrack, Application web et mobile, Microservices, Intelligence Artificielle, RAG, Géolocalisation, Géofencing, Sécurité proactive, Flutter, NestJS, MongoDB.

\vspace{1em}

\begin{center}
\textbf{\large Abstract : }
\end{center}

This report presents the development of \textbf{SafeTrack}, an intelligent web and mobile application designed for real-time safety monitoring and tracking of children and vulnerable individuals. Built on a microservices architecture, the platform integrates several features including secure authentication, real-time GPS tracking, geofencing management, SOS alerts, instant notifications, and AI-powered predictive risk analysis.

The system leverages advanced artificial intelligence techniques, particularly \textbf{Retrieval-Augmented Generation (RAG)}, to provide contextual safety analysis and personalized recommendations based on detected situations. The application also enables real-time communication between parents, caregivers, and institutions to enhance proactive protection.

The main technologies used include \textbf{Flutter}, \textbf{NestJS}, \textbf{Python}, \textbf{MongoDB}, \textbf{Redis}, \textbf{WebSocket}, along with AI models integrated through \textbf{Ollama} and \textbf{LangChain}.\\

\noindent
\textbf{Keywords:} SafeTrack, Web and Mobile Application, Microservices, Artificial Intelligence, RAG, Geolocation, Geofencing, Proactive Safety, Flutter, NestJS, MongoDB.